\section{Фано-таблицы: фабрика, медиаторы, синдром}
\label{app:fano}

\begin{sloppypar}
Все таблицы машинно верифицированы (\texttt{domain\_transfer\_verify.py},
\texttt{talos\_collapse\_uniqueness.py}).
\end{sloppypar}

\paragraph{Семь линие-шин.} Фабрика \S\ref{sec:fabric}: $\{A,S,L\}$, $\{A,D,E\}$,
$\{A,O,U\}$, $\{S,D,O\}$, $\{S,E,U\}$, $\{D,L,U\}$, $\{L,E,O\}$. Всякий сектор
лежит ровно на трёх; всякая пара --- ровно на одной.

\paragraph{Карта медиаторов.} Бестабличная перемаршрутизация \S\ref{sec:fabric}:
для всякой пары $(i,j)$ медиатор $k^\ast(i,j)$ --- третий сектор их линии, и
единственный путь, которым регенерирует их когерентность.

\begin{table}[H]
\centering\small
\begin{tabular}{@{}llll llll@{}}
\toprule
пара & $k^\ast$ & пара & $k^\ast$ & пара & $k^\ast$ & пара & $k^\ast$ \\
\midrule
$AS$ & $L$ & $AD$ & $E$ & $AL$ & $S$ & $AE$ & $D$ \\
$AO$ & $U$ & $AU$ & $O$ & $SD$ & $O$ & $SL$ & $A$ \\
$SE$ & $U$ & $SO$ & $D$ & $SU$ & $E$ & $DL$ & $U$ \\
$DE$ & $A$ & $DO$ & $S$ & $DU$ & $L$ & $LE$ & $O$ \\
$LO$ & $E$ & $LU$ & $D$ & $EO$ & $L$ & $EU$ & $S$ \\
$OU$ & $A$ &      &     &      &     &      &     \\
\bottomrule
\end{tabular}
\caption{$21$ медиатор. Читать: чтобы починить или маршрутизировать $(i,j)$,
используй плечи через $k^\ast(i,j)$. Нет таблицы маршрутизации --- $k^\ast$
вычисляется из пары.}
\end{table}

\paragraph{Декодирование синдрома.} Локализация отказа при нулевых накладных
\S\ref{sec:ecc}. Пометим семь секторов $1\dots7$ ненулевыми векторами
$\mathbb{F}_2^3$. Единичный отказ в секторе $p$ производит трёхбитовый синдром,
равный двоичному $p$; три сработавшие проверки чётности --- три линии через $p$.
Двадцать одно парное отношение коллапсирует к семи осям, читаемым тремя битами
синдрома, к одному вердикту --- пирамида $21\to7\to3\to1$.

\begin{table}[H]
\centering\small
\begin{tabular}{@{}ccl@{}}
\toprule
сектор отказа $p$ & синдром (двоичн.) & линии через $p$ (3 сработавшие проверки) \\
\midrule
1 & $001$ & три линии, содержащие сектор 1 \\
2 & $010$ & три линии, содержащие сектор 2 \\
3 & $011$ & три линии, содержащие сектор 3 \\
4 & $100$ & три линии, содержащие сектор 4 \\
5 & $101$ & три линии, содержащие сектор 5 \\
6 & $110$ & три линии, содержащие сектор 6 \\
7 & $111$ & три линии, содержащие сектор 7 \\
\bottomrule
\end{tabular}
\caption{Хэмминг $[7,4,3]$, синдром единичной ошибки: синдром, читаемый как
двоичный индекс, \emph{есть} позиция отказа (проверено для всех семи).
Декодирование --- табличный поиск глубины три; нет выделенного хранения проверок.}
\end{table}

\paragraph{Арифметика уникальности.} Теорема~\ref{thm:unique}: $q^2+q+1=2^r-1$
имеет целые решения $q\in\{2,5,90\}\Rightarrow N\in\{7,31,8191\}$, из которых
лишь $q=2,5$ (т.\,е.\ $N=7,31$) суть степени простого и потому подлинные
проективные порядки ($q=90$ не степень простого). Мнимые размерности алгебр с
делением суть $\{0,1,3,7\}$; неассоциативная --- лишь $7$; тройное пересечение ---
$\{7\}$.
